\graphicspath{{abstracts/praesentationen/05-linked-open-data/img/}{img/}}

\begin{beitrag}{Linked Open Data im musealen Kontext: Ein Erfahrungsbericht aus fünf Jahren Sammlungsdigitalisierung}{%
Iris Ihlow\,\orcidlink{0000-0003-0123-4567} \\
\small \href{mailto:iris.ihlow@example.org}{iris.ihlow@example.org} \\
\small Universität Hamburg \\[0.5cm]
Jan Jansen\,\orcidlink{0000-0001-1234-5678} \\
\small \href{mailto:jan.jansen@example.org}{jan.jansen@example.org} \\
\small Rheinische Friedrich-Wilhelms-Universität Bonn
}

\section{Einleitung}
Die Idee, museale Sammlungsdaten als Linked Open Data (LOD) zu publizieren, ist seit weit über einem Jahrzehnt fester Bestandteil der einschlägigen Diskurse. In der praktischen Umsetzung zeigt sich jedoch immer wieder, dass die methodische Strenge offener semantischer Datenmodelle nur schwer mit den real existierenden Sammlungsdatenbanken, ihren historischen Klassifikationssystemen und ihren institutionellen Rahmenbedingungen vereinbar ist. Der vorliegende Beitrag zieht eine Zwischenbilanz nach fünf Jahren LOD-Praxis an einem mittelgroßen kulturhistorischen Museum und reflektiert über das, was funktioniert hat -- und über das, was wir heute anders machen würden.

\section{Was funktioniert hat}
Die zentrale Stärke der LOD-Strategie liegt in ihrer disziplinierenden Wirkung auf die Datenerfassung. Schon der Versuch, eine Klassifikation \glqq Möbel\grqq{} in eine kontrollierte Vokabular-URI zu überführen, hat im Projektteam zu einer fundierten Diskussion über Begriffshierarchien, regionale Unterschiede und die Frage geführt, ob die Sammlung sich primär als regionalhistorisch oder als typologisch versteht. Diese Diskussionen waren produktiv und haben langfristig zu einer schärferen Profilbildung der Sammlung beigetragen.

Technisch hat sich der Workflow bewährt, in dem WissKI als zentrale Datenerfassungs- und Erschließungsumgebung eingesetzt wird, während die LOD-Publikation über einen vorgelagerten Aufbereitungsschritt erfolgt, der die kuratierte Erschließung in stabile URIs überführt und mit Wikidata, GND und Getty AAT abgleicht. Eine ausführliche Darstellung des Workflows findet sich bei \cite{ihlow2024-museum-lod}.

\section{Was schwieriger war als erwartet}
Drei Themenfelder haben sich als deutlich anspruchsvoller erwiesen, als wir zu Beginn des Projekts angenommen haben. Erstens die Frage des \emph{Granularitätsausgleichs}: Wenn ein internes Klassifikationssystem etwa 1\,400 Begriffe umfasst und die Bezugs-Ontologie nur 380, müssen Mappings entwickelt werden, die Verluste nicht verschleiern. Wir haben dafür ein dreigliedriges Mapping-Schema eingeführt (\texttt{skos:exactMatch}, \texttt{skos:closeMatch}, \texttt{skos:broadMatch}), das im Frontend für die fachliche Auseinandersetzung sichtbar gemacht wird.

Zweitens die \emph{Stabilität externer Bezugsdaten}: Wikidata-Identifikatoren werden gelegentlich umgewidmet, Getty-Konzepte mit der Zeit umgehängt, GND-Datensätze zusammengelegt. Jede dieser Veränderungen kann das eigene Datenmodell erschüttern, wenn sie nicht systematisch beobachtet wird. Wir haben deshalb einen wöchentlichen Konsistenzcheck eingeführt, der gegen die jeweils kanonischen Endpunkte arbeitet und auffällige Veränderungen meldet.

Drittens die \emph{institutionelle Verstetigung}: Der laufende Betrieb einer LOD-Publikation erfordert eine kontinuierliche personelle Bereitstellung, die in einer projektbezogenen Drittmittelförderung schwer abbildbar ist. Wir plädieren in diesem Zusammenhang dafür, LOD-Pflege explizit als Dauerlasten der Museums-IT zu verankern -- analog zur Pflege bibliografischer Datenbanken; siehe \cite{jansen2025-lod-finanzierung}.

\section{Reflexion und Empfehlungen}
Aus fünf Jahren Praxis ergibt sich für uns folgendes Bild: LOD ist kein Selbstzweck und nur dann sinnvoll, wenn die institutionellen Voraussetzungen für eine dauerhafte Pflege gegeben sind. Wo sie gegeben sind, lohnt sich der Aufwand mehrfach -- die eigene Sammlung wird besser erschlossen, ihre Sichtbarkeit steigt, und die Anschlussfähigkeit an Forschungsprojekte verbessert sich substantiell. Wo die Voraussetzungen nicht gegeben sind, sollte zunächst in die interne Datenqualität investiert werden, bevor eine vorzeitige Publikation den Eindruck von Solidität erzeugt, dem die Datenbasis nicht entspricht.

\section{Ausblick}
Im laufenden Jahr werden wir den Schritt zur Anbindung an das europäische Datenraum-Netzwerk \emph{CulturalData.eu} unternehmen und freuen uns auf den Erfahrungsaustausch mit Einrichtungen, die einen ähnlichen Weg planen.

\end{beitrag}
